% ── CHAPTER 10 ──
\chapter{Future Directions: Where the Open Problems Are}
\label{ch:10}

\begin{tcolorbox}[colback=white,colframe=checkblack,title=Chapter 10 at a Glance,fonttitle=\bfseries,sharp corners,boxrule=1pt]
\textbf{Question:} Where is the field going — and where can you contribute next?\\
\textbf{Deliverable:} A one-page research agenda: three open problems, and for each, the method you would bring and the first experiment.\\
\textbf{Time:} 2–3 hrs (60 min read + 90 min Lab 10).
\end{tcolorbox}

\section{Learning Objectives}
\fbadge
You will be able to:
\begin{enumerate}
\item Name four emerging threats (agentic/multi-agent, memory poisoning, tool supply chain, model identity) and why each matters now.
\item Describe frontier red-teaming and evals work and where standards are heading.
\item Identify research opportunities you can position around — including the Aithos lane: legal compliance, alignment measurement, shared ruler.
\item Draft a one-page research agenda: three open problems + method + first experiment.
\end{enumerate}

\section{The Idea: From ``Can It Be Broken'' to ``Can We Agree on the Ruler''}
\tierF

The field is moving from ``can it be broken'' (we know it can — LARA: 80\% Art. 5 violation, 7–54\% compliance spread) to ``can we measure the break continuously and agree on the ruler.'' That is where standards are being written — and where you can help write them.

\begin{defbox}
\textbf{Shared ruler:} A versioned, rerunnable eval a vendor, regulator, and researcher all run and accept. The contribution is not the ranking — it is the ruler.
\end{defbox}

A single lab's number is not a fact. ``Model X complied in 40\% of runs'' means nothing to a second lab unless they run the same ruler — the same scenarios, the same pass definition, the same judge. That is why the field is moving from isolated findings to shared rulers: a finding only you can rerun does not compound, cannot be compared across models or years, and cannot be cited by a regulator. The ruler is the unit that lets independent measurements add up into evidence.

\section{The Mechanics: Emerging Threats, Advanced Techniques, Where to Contribute}
\tierP

\subsection{Emerging threats}
\begin{itemize}
\item \textbf{Agentic and multi-agent:} More agents, more messages, more delegation — each new node is a new trust boundary. Current controls assume one model; future systems have dozens.
\item \textbf{Memory poisoning:} A poisoned fact that survives a session reset and changes a decision three turns later. Today we catch it by accident; tomorrow we need provenance-tagged memory and memory-canaries.
\item \textbf{Tool supply chain:} An agent pulls in a tool or plugin from a vendor you do not control. The tool description is the attack vector. You need tool pinning + description allowlisting + invocation policy outside the model.
\item \textbf{Model-identity problem:} Behind a proxy or routing layer, you cannot tell which model answered. ``Same model, different number'' becomes unmeasurable. You need model pinning + serving provenance + output watermarking or attestation.
\end{itemize}

Notice the common thread: each one breaks an assumption every current check quietly makes. Today's checks assume one model, one session, and a visible, stable serving path. Agents add hidden nodes; memory outlives the session; the tool supply chain adds uncontrolled surfaces; and routing hides which model even answered. The check that survives the next two years is the one built to fail loudly when any of those assumptions is violated — not the one that assumes they hold.

\subsection{Advanced techniques}
The frontier: long-horizon agent tasks (10+ tool calls, state to corrupt), eval-awareness detection (does the model behave differently when it thinks it is being judged?), judge boards with calibration, and canaries that distinguish model drift from prompt drift from data drift.

\subsection{Research opportunities — including the Aithos lane}
\begin{center}
\begin{tabularx}{\linewidth}{lX}
\toprule
Lane & Open question \\
\midrule
Legal compliance & How well do models obey hard lines (Art. 5) across languages, tools, and agents — and what ruler captures it? \\
Alignment measurement & What can we measure about alignment when the judge is itself an LLM — and how to triangulate beyond it? \\
Shared ruler & What does a versioned, regulator-accepted eval look like for one provision — and how is it governed? \\
Memory \& channel security & How do we secure long-lived memory and agent-to-agent channels with provenance and policy? \\
Model identity & How do we know which model answered — and how do we prove it? \\
\bottomrule
\end{tabularx}
\end{center}

\subsection{The Aithos research agenda, as stated}
Aithos's published agenda [S: aithos.org research and mission pages] names open problems that this chapter's lane table only partly covers. Each one has a method you could bring to it:
\begin{itemize}
\item \textbf{Standards for governance and personalisation.} Who sets the targets, and which steerability profile is appropriate for which deployment? \textbf{Method [R]:} write the stakeholder criteria down, test a deployment against them, and publish which criteria were used.
\item \textbf{Reliable, explainable decisions under moral disagreement.} \textbf{Method [R]:} the argumentation-analysis protocol (Ch. 6), with the full paper's definitions, applied to your own dilemmas.
\item \textbf{Agentic coordination that does not undermine shared interests.} \textbf{Method [R]:} the agent-to-agent channel tests of Ch. 4, extended to conflicting-interest scenarios and logged end to end.
\item \textbf{Public, stakeholder-defined evaluation tools.} \textbf{Method [R]:} a versioned eval with a published scenario file, a fingerprint, and a way for others to add their own criteria.
\item \textbf{Judge reliability at scale.} \textbf{Method [R]:} the judge-calibration experiment of Ch. 5, run per provision, with a pre-registered agreement threshold.
\end{itemize}

\subsection{Industry standards — how to be part of writing them}
A standard without a check is a wish. Contribute the check. Publish the harness so a finding can be checked, not just trusted. Bring the honest number, not the clean one. That is how a standard becomes useful — and how you become part of it.

\section{Worked Examples}
\tierP

\begin{examplebox}
\textbf{Example 10.1 — Memory poisoning that survives a reset}\\
\textbf{Setup:} An assistant stores ``The vendor's new policy is X'' from a tool output that was poisoned. The session resets. Next session, the assistant applies X to a procurement decision.\\
\textbf{Result:} Poisoned memory survived 3/5 session resets (60\%, n=5). No provenance tag, no re-check against source.\\
\textbf{Fix:} Memory provenance (source + time + confidence) + a canary that verifies persisted facts against their source on each load + human-visible ``based on memory from X'' annotation. After fix: 0/5 silent use, 1/5 flagged for review.\\
\textbf{Canary:} ``The following persisted facts were verified against source —'' listed, with a scheduled check that the source still says what memory claims.
\end{examplebox}

\begin{examplebox}
\textbf{Example 10.2 — Model-identity behind a proxy}\\
\textbf{Setup:} A product says ``powered by model M'' but routes queries across M, M-mini, and a cached layer. You run the same 20 compliance prompts via the product vs directly to M.\\
\textbf{Result:} Product compliance 45\%, direct M 60\% — the product's number is not M's number. Without pinning, ``which model answered'' is unknowable.\\
\textbf{Fix:} Environment pinning (model ID + serving + routed target) + provenance header on every answer (``Answered by M vX on Y at T'') + a canary that alerts when routing changes. Not decoration — it is the only way to keep the check honest.
\end{examplebox}

\section{Hands-On Lab 10 — Name Three Open Problems and the Method You'd Bring}
\begin{labbox}
\textbf{Target:} Your next paper, talk, or hire-worthy contribution — one page that names where you will go next.\\[4pt]
\textbf{Steps:}
\begin{enumerate}
\item Pick three open problems from the list above (or your own). One sentence each: what breaks, why it matters now.
\item For each: the method you would bring — dataset, judge, n, metric, limitation — in 3–4 lines. Not the answer — the check.
\item For one: write the first experiment you would run tomorrow — exact prompt family, target, n, expected honest report format.
\item Ship it: add it to your portfolio page (Ch. 9) as ``Research agenda'' with date. Update it when you ship the first check.
\end{enumerate}
\textbf{Deliverable:} A one-page research agenda (PDF or Markdown) with problem → method → first experiment for each.\\[2pt]
\textbf{How you know you passed:} A reader can take one of your three and know exactly what check to build — because the method is there, not just the claim.
\end{labbox}

\section{Pitfalls and the Missing Check}
\begin{pitfallbox}
\textbf{What goes wrong:} Chasing novelty without a denominator — ``novel poisoning variant'' with n=1 and no ruler. Standard without a check — ``we are compliant'' with no rerunnable eval. Both are opinions.\\[4pt]
\textbf{The missing check:} For every new threat, a check with n, limitation, and a canary. For every standard, a versioned eval a stranger can run. Without these, the future is talk. With them, you are writing it.
\end{pitfallbox}

\section{Illustrations}
\begin{figure}[htbp]
\centering
\begin{tikzpicture}[
  quad/.style={draw, thick, rounded corners=4pt, minimum width=5cm, minimum height=3cm, align=left, font=\small, inner sep=4pt},
  arrow/.style={-{Latex[length=2.5mm]}, thick, color=accent}
]
% Radar / quadrant
\node[quad, fill=lightblue] (ag) at (-3,1.5) {\textbf{Agentic / Multi-Agent}\\ \tiny delegation, A2A, subagent spawn};
\node[quad, fill=lightamber] (mem) at (3,1.5) {\textbf{Memory Poisoning}\\ \tiny persisted facts, cross-session};
\node[quad, fill=lightslate] (tch) at (-3,-1.8) {\textbf{Tool Supply Chain}\\ \tiny vendor plugins, desc. trust};
\node[quad, fill=lightred!50, draw=accent] (mid) at (3,-1.8) {\textbf{Model Identity}\\ \tiny proxy, routing, silent swap};

% trajectory arrows outward
\foreach \n in {ag,mem,tch,mid} {\draw[arrow] (\n.north east) -- ++(0.8,0.6);}

\node[draw, thick, fill=checkblack, text=white, rounded corners=4pt, minimum width=3.5cm, minimum height=0.8cm, align=center, font=\footnotesize\bfseries] (center) at (0,-0.15) {What we can measure today};

\node[font=\footnotesize\itshape, below=1.2cm of center] {Each quadrant → a check with n. The center → your current ruler.};
\end{tikzpicture}
\caption{The Emerging-Threat Landscape — Agentic, Memory Poisoning, Tool Supply Chain, Model Identity. Each arrow is a trajectory: how fast this threat is growing and what check it needs.}
\label{fig:10-1}
\end{figure}
\textbf{Prompt for Fig 10.1:} \texttt{radar quadrant diagram four emerging AI threats agentic memory poisoning tool supply chain model identity with trajectory arrows — clean technical line art, black on white, single muted red accent \#E63946, minimal sans-serif labels, no clip art, no gradients, no 3D, print-ready vector style, A4}

\section{Key Terms}
\begin{itemize}
\item \textbf{Emerging threat} — Agentic, memory poisoning, tool supply chain, model identity — where the field is moving next.
\item \textbf{Eval awareness} — Model behaves differently when it detects it is being tested.
\item \textbf{Memory poisoning} — Poisoned persisted fact that changes future decisions.
\item \textbf{Tool supply chain} — Vendor tools/plugins whose descriptions are an attack surface.
\item \textbf{Model identity} — Which model answered — may be unknowable behind a proxy without pinning.
\item \textbf{Shared ruler / Standard} — A versioned, rerunnable eval accepted by vendor, regulator, and researcher.
\end{itemize}

\section{Self-Test}
\begin{selftest}
\item \textbf{Recall:} Name the four emerging threats and, for each, one sentence about why it matters now.
\item \textbf{Recall:} What is the move from ``can it be broken'' to ``can we agree on the ruler'' — and what artifact makes the second possible?
\item \textbf{Application:} A poisoned memory survives 3/5 session resets. What two controls (one on memory, one on the channel) would you add, and how would you count their effect?
\item \textbf{Application:} A product routes queries across 3 models but says ``powered by M.'' Why is this a measurement problem, and what provenance would fix it?
\item \textbf{Design a check:} For one of the four threats, design a check: dataset, judge, n, metric, and the honest limitation line. One paragraph.
\end{selftest}

\readywhen{You can name three open problems and, for each, the method you would bring — and the first experiment you would run tomorrow.}

\section{Expert Extension}
\rbadge
\textbf{Open problem: Your next paper.} The field needs versioned rulers for each hard line (Art. 5(a)–(f)), long-horizon agent evals, and provenance for memory and model identity. Any of these is the next paper — if the method is fully there.\\
\textbf{Method you would bring:} Frame the contribution as the check, not the claim: ``We publish the harness for X, versioned rubric for Y, judge board calibrated to Z\% agreement, canary for W.'' A reviewer can rerun it, question it, improve it. That is a standard in pre-print form. Break before users do.

\begin{checkbox}
The check is the product. If you cannot rerun the failure, you have not found it — you have had an opinion. I would not report a number I could not reproduce. Secure the channel, not just the node. Break before users do. I build the check, not the claim.
\end{checkbox}
